
\documentclass{article}
\usepackage{colortbl}
\usepackage{makecell}
\usepackage{multirow}
\usepackage{supertabular}

\begin{document}

\newcounter{utterance}

\centering \large Interaction Transcript for game `cladder', experiment `full\_v1.5\_default', episode 6916 with Qwen3.5{-}9B.
\vspace{24pt}

{ \footnotesize  \setcounter{utterance}{1}
\setlength{\tabcolsep}{0pt}
\begin{supertabular}{c@{$\;$}|p{.15\linewidth}@{}p{.15\linewidth}p{.15\linewidth}p{.15\linewidth}p{.15\linewidth}p{.15\linewidth}}
    \# & \multicolumn{2}{c}{Player} && \multicolumn{2}{c}{Game Master} \\
    \hline

    \theutterance \stepcounter{utterance}  
    & & & \multicolumn{4}{p{0.6\linewidth}}{
        \cellcolor[rgb]{0.9,0.9,0.9}{
            \makecell[{{p{\linewidth}}}]{
                \texttt{\tiny{[P1$\langle$GM]}}
                \texttt{You are an expert in causal inference. The following question is not a typical commonsense query, but rather a meticulously designed question created by a professor specializing in causal inference, intended to assess the students' mastery of the course content.} \\
\texttt{Imagine a self{-}contained, hypothetical world with only the following conditions, and without any unmentioned factors or causal relationships: Jyka has a direct effect on hwax. Lirg has a direct effect on hwax. Method 1: We look at how jyka correlates with lirg case by case according to hwax. Method 2: We look directly at how jyka correlates with lirg in general. To understand how jyka affects lirg, is it more correct to use the Method 1 than Method 2?} \\
\texttt{Start your answer with 'Yes' or 'No', followed by additional reasoning or evidence to support your explanation.} \\
            }
        }
    }
     \\ \\

    \theutterance \stepcounter{utterance}  
    & \multicolumn{4}{p{0.6\linewidth}}{
        \cellcolor[rgb]{0.561,0.737,0.561}{
            \makecell[{{p{\linewidth}}}]{
                \texttt{\tiny{[P1$\rangle$GM]}}
                \texttt{No} \\
\\ 
\texttt{In this hypothetical world, **Jyka** and **Lirg** are both direct causes of the same effect, **hwax**. This structure creates a classic **confounding** scenario (specifically, a "fork" or common cause structure if we view hwax as the outcome, but here we are looking at the relationship between the two causes).} \\
\\ 
\texttt{Here is the causal reasoning:} \\
\\ 
\texttt{1.  **The Causal Structure**:} \\
\texttt{    *   \$Jyka \textbackslash{}rightarrow hwax\$} \\
\texttt{    *   \$Lirg \textbackslash{}rightarrow hwax\$} \\
\texttt{    *   There is no direct arrow between Jyka and Lirg.} \\
\\ 
\texttt{2.  **Analyzing Method 2 (Unconditional Correlation)**:} \\
\texttt{    *   Method 2 looks at the raw correlation between Jyka and Lirg without conditioning on anything.} \\
\texttt{    *   In a purely hypothetical world with *only* the stated conditions, if Jyka and Lirg are independent variables (no unmentioned factors linking them), their unconditional correlation would be zero.} \\
\texttt{    *   However, if there is any latent selection bias or if the variables are not independent in the real{-}world data generation process, Method 2 captures the **total association**, which includes any spurious correlation induced by the fact that they both influence the same outcome (though typically, correlation between causes is not induced by the effect unless there is feedback or selection on the outcome).} \\
\texttt{    *   Crucially, Method 2 does not isolate the direct causal link because there *is* no direct link. It simply measures the marginal association, which is likely zero or driven by external, unmentioned factors.} \\
\\ 
\texttt{3.  **Analyzing Method 1 (Conditional Correlation)**:} \\
\texttt{    *   Method 1 looks at the correlation between Jyka and Lirg **case by case according to hwax**. This phrasing implies conditioning on the value of **hwax** (i.e., calculating the correlation within strata of hwax, or performing a partial correlation controlling for hwax).} \\
\texttt{    *   In causal inference, when two variables share a common effect (a "collider"), conditioning on that common effect opens a non{-}causal path between them. This is known as **collider bias** or **M{-}bias** (depending on the specific graph orientation, but here it is a fork where we condition on the node).} \\
\texttt{    *   Mathematically, if \$Jyka \textbackslash{}perp Lirg\$ (independent), but both cause \$hwax\$, then \$Jyka \textbackslash{}not\textbackslash{}perp Lirg \textbackslash{}mid hwax\$. Conditioning on \$hwax\$ induces a spurious negative correlation between Jyka and Lirg.} \\
\texttt{    *   Therefore, Method 1 would reveal a **spurious correlation** that does not exist in the causal graph. It creates a false impression that Jyka affects Lirg (or vice versa)} \\
            }
        }
    }
    & & \\ \\

    \theutterance \stepcounter{utterance}  
    & & & \multicolumn{2}{p{0.3\linewidth}}{
        \cellcolor[rgb]{0.9,0.9,0.9}{
            \makecell[{{p{\linewidth}}}]{
                \texttt{\tiny{[GM$|$GM]}}
                \texttt{No} \\
            }
        }
    }
    & & \\ \\

    \theutterance \stepcounter{utterance}  
    & & & \multicolumn{2}{p{0.3\linewidth}}{
        \cellcolor[rgb]{0.9,0.9,0.9}{
            \makecell[{{p{\linewidth}}}]{
                \texttt{\tiny{[GM$|$GM]}}
                \texttt{no} \\
            }
        }
    }
    & & \\ \\

    \theutterance \stepcounter{utterance}  
    & & & \multicolumn{2}{p{0.3\linewidth}}{
        \cellcolor[rgb]{0.9,0.9,0.9}{
            \makecell[{{p{\linewidth}}}]{
                \texttt{\tiny{[GM$|$GM]}}
                \texttt{game\_result = WIN} \\
            }
        }
    }
    & & \\ \\

\end{supertabular}
}

\end{document}
